% Additive reconstruction of Untitled 1775.md, equations (9)--(20),(25).
% Inclusion file: requires amsmath,amssymb,hyperref.
\begingroup
\section{The corrected native radial profile and its endpoint defect}
\label{sec:CRS}

This calculation receives the literal conductor-pivot identity from the
accompanying correlated-row proof and proves the scalar transformations,
finite error transfers and coefficients. Its source is the first
continuation in \texttt{Untitled 1775.md}, equations (9)--(20) and (25).
Earlier providers are LRS1--18 and LRS24--25, PCB1--6, the direct elliptic
derivative proof and ATA35--40. Their exact files are pinned in the
adjacent \texttt{ACCEPTANCE.json}. No sealed predecessor is edited.

\subsection{Original parameters and finite input identities}

Keep the actual simple quartet, period, branch and conductor. Its moment
is $\mu_v=E_A^{(v)}(0)\ne0$, $0\le v\le80$. Write
\[
\begin{gathered}
 a\equiv1\pmod4,\quad a\ge257,\quad s\in\{1,a\},\quad m=\Delta-v,\\
 q=(a+1)^2,\quad Q=(a-7)^2,\quad \Delta=q-Q=16a-48,\\
 \ell=(s-1)/4,\quad n_+=q+\ell,\quad n_-=Q+\ell,\quad
 \alpha_A=v!/\mu_v,\quad \ell_r=\mu_v\binom{q+r}{v}.
\end{gathered}\tag{CRS1}
\]
The degree-$a-8$ lower packet keeps the same order $s$. Retain the two
original LRC13 guards:
\[
 n_b=(b+1)^2+\ell\ge100,\qquad
 \max\{b\sqrt{\delta^2+\gamma^2},2\ell+1/2\}
 \le2^{-17}n_b,\qquad b=a,a-8.                         \tag{CRS2}
\]
Any GPA provider used elsewhere keeps its own guard. The present
calculation uses the retained Gamma LRC endpoints with their original
weight. All sums allow $0\le M\le q+32a$.
The original degree ranges hold: $M\le2n_+$ and $m+M\le2n_-$.
For the latter,
$q-3\Delta-32a+v+2\ell=a^2-78a+145+v+2\ell>0$ for $a\ge257$.
All scalar arguments below lie in $[0,3]$, since
$M+\Delta\le q+48a-48<3Q$.

Let $\nu^+_{r,s},\nu^-_{m+r,s}$ be the original full monic minima and
$\eta_{r,s}$ the attained preceding-relation minimum. Put
\[
 h_{d,s}=(2\pi)^{s/2}d!(s/2)_d,\quad
 \gamma_{r,s}=\log(\eta_{r,s}/\nu^-_{m+r,s})\ge0.
\]
The squared Cholesky pivot $\delta_r$ of the complete conductor-image
Gram in its original degree order satisfies
\[
 \delta_r=|\ell_r|^2\eta_{r,s}/\nu^+_{r,s}.             \tag{CRS3}
\]
This is the affine fibre identity ATA6 followed by the full Schur
minimum, as reconstructed in the accompanying conductor proof.

For root count $p=q,Q$, let $e_{p,j,s}$ be the larger absolute distance
of the two finite LRS4 endpoints from $2(p+\ell)\psi(j/(p+\ell))$.
Use the scalar-moment endpoints at $j=0,1$. LRS11 proves
\[
\begin{split}
 \log\nu^{(p)}_{j,s}-\log h_{p+j,s}
 &=2n\psi(j/n)+\epsilon_{p,j,s},\quad
 |\epsilon_{p,j,s}|\le e_{p,j,s},\quad n=p+\ell,\\
 e_{p,j,s}&\le C_{\rm actual}
 \left(1+\log^+\frac{n}{j+1}+\frac{(\ell+1)^2}{n+j}\right),
 \qquad 0\le j\le2n .
\end{split}                                                    \tag{CRS4}
\]
Its retained Gamma parity parameter is $(n+\varepsilon-3/4)/h$.
The positive-real Stirling remainder used in LRS7 has human provenance
in NIST DLMF 5.11.1 and 5.11(ii) \cite[\S5.11]{crs:dlmf5}.
The norm masses and parity parameter are unchanged here.
The errors are summable: integrating $\log(n/x)$ on $(0,n)$ gives
$\sum_{j=0}^{2n}\log^+(n/(j+1))\le n$, and
$\sum_{j=0}^{2n}(\ell+1)^2/(n+j)\le3(\ell+1)^2=O(q)$.
There are $O(q)$ constant terms. Thus both original row-error sums
have size $O_{\rm actual}(q)$.

\subsection{The exact elliptic dictionary}

Use the real Legendre modulus convention in NIST DLMF 19.2.4,
19.2.5 and 19.2.8 \cite[\S19.2]{crs:dlmf19}:
\[
 K(k)=\int_0^{\pi/2}(1-k^2\sin^2\theta)^{-1/2}\,d\theta,\quad
 E(k)=\int_0^{\pi/2}(1-k^2\sin^2\theta)^{1/2}\,d\theta.
\]
The original parameter and functions are
\[
\begin{gathered}
 k(t)=\sqrt{1-\rho(t)^2},\qquad
 \rho(t)K(k(t))/E(k(t))=(1+t)^{-1},\\
 f(t)=\log\frac{1+\rho(t)}{1-\rho(t)},\qquad
 \psi(t)=\log\frac{1+\rho(t)}{E(k(t))}
                         +t\log\frac{k(t)}{E(k(t))}.
\end{gathered}                                                   \tag{CRS5}
\]
Equivalently $\psi(t)=(1+t)(1-\log(1+t))-t\lambda(2/t,\pi/t)/4$
with the original Robin convention. The pinned direct integral proof
establishes uniqueness of $\rho$, the derivative identities, and
the following global and endpoint bounds:
\[
\begin{gathered}
 \psi'(t)=\log(k(t)/E(k(t)))<0,\quad
 2\psi(t)-2(1+t)\psi'(t)=f(t),\quad \psi(0)=a_0=\log(4/\pi),\\
 0<-tf'(t)\le1,\quad
 |f(t)+\tfrac12\log t|\le(22/5)\sqrt t,\quad
 |tf'(t)+\tfrac12|\le9\sqrt t\quad(0<t\le1/20).
\end{gathered}                                                   \tag{CRS6}
\]
For the dictionary identity, substitute CRS5: its left side is
$2\log((1+\rho)/k)=\log((1+\rho)/(1-\rho))$ exactly.
The derivative bounds in the provider are proved by differentiating
the displayed defining integrals, without differentiating an
asymptotic remainder.
Put
\[
 \eta(t)=f(t)+\tfrac12\log t,\quad \eta(0)=0,\qquad
 H_\eta=\int_0^3|\eta'(t)|\,dt
          \le18/\sqrt{20}+\tfrac12\log60.              \tag{CRS7}
\]
Indeed $|\eta'|\le9/\sqrt t$ up to $1/20$, and thereafter
$-tf'\in[0,1]$ gives $|\eta'|\le1/(2t)$. This proves absolute
continuity of $\eta$ on $[0,3]$ and the finite bound.
Also $\int_0^3|\psi'|=\psi(0)-\psi(3)<\infty$.

\subsection{A finite summed scalar error}

Define
\[
 h_r^{\rm prof}=2n_+\psi(r/n_+)-2n_-\psi((m+r)/n_-)+\Gamma_r.
                                                               \tag{CRS8}
\]
The factorial correction is exactly
\[
\begin{split}
 \Gamma_r&=\log\frac{h_{q+r,s}}{h_{q+r-v,s}}-2\log|\ell_r|\\
 &=2\log|\alpha_A|+
   \sum_{j=0}^{v-1}\log\left(1+\frac{s/2-1}{q+r-j}\right).
\end{split}                                                    \tag{CRS9}
\]
The norm ratio contributes
$\prod_{j<v}(q+r-j)(q+r-j-1+s/2)$ and
$|\ell_r|^2=|\mu_v/v!|^2\prod_{j<v}(q+r-j)^2$.
Their quotient proves CRS9 with the original mass cancelled only
inside this matched ratio. Empty products apply at $v=0$.
Combining CRS3--4 gives the exact equality
\[
 \gamma_{r,s}-h_r^{\rm prof}
   =\log\delta_r+\epsilon_{q,r,s}-\epsilon_{Q,m+r,s}.  \tag{CRS10}
\]
Let $B_M=2F_*+2I_*+C_{\rm ang}$ be the finite bound for
$\sum_{r=0}^M|\log\delta_r|$ proved in the accompanying conductor
calculation. Its proved size is $O_{\rm actual}(q)$.
Consequently
\[
 E_1=B_M+\sum_{r=0}^M(e_{q,r,s}+e_{Q,m+r,s}),\qquad
 \sum_{r=0}^M|\gamma_{r,s}-h_r^{\rm prof}|
 \le E_1=O_{\rm actual}(q).                           \tag{CRS11}
\]
No operator-norm error is charged independently to each pivot.

Set $F(n,r)=2n\psi(r/n)$. Direct differentiation and CRS6 give
$F_n-F_r=f(r/n)$. Integrate on $(n,r)=(n_+-u,r+u)$ and retain
the actual $m=\Delta-v$:
\[
\begin{split}
 F(n_+,r)-F(n_-,r+m)
 =\int_0^\Delta f\left(\frac{r+u}{n_+-u}\right)\,du
       +2\int_{r+m}^{r+\Delta}\psi'(b/n_-)\,db .
\end{split}                                                    \tag{CRS12}
\]
At $r=0$ these are convergent improper integrals by CRS6.
For $x>0$ put
\[
 c(x)=\Delta f(x/q)
 -\tfrac12\{(x+\Delta)\log(1+\Delta/x)-\Delta\}.        \tag{CRS13}
\]
Its exact integral presentation, valid also at $x=0$, is
\[
 c(x)=-\tfrac12\int_0^\Delta\log((x+u)/q)\,du
                     +\Delta\eta(x/q),\qquad
 c_0=c(0)=\tfrac{\Delta}{2}(1+\log(q/\Delta)),\quad c_r=c(r).
                                                               \tag{CRS14}
\]
This proves the initial-row constant with no unspecified offset.

All scalar terms between CRS12 and CRS13 have the following
finite summed bound:
\[
\begin{split}
 B_\Gamma&=(M+1)
 \left(2|\log|\alpha_A||+\frac{2v|s/2-1|}{q-v+1}\right),\\
 B_v&=2(v+1)n_-[\psi(0)-\psi(3)],\\
 B_{\rm den}&=(M+1)\Delta(\ell+\Delta)/Q,\qquad
 B_\eta=(\Delta^2/2+\Delta)H_\eta,\\
 B_{\rm sc}&=B_\Gamma+B_v+B_{\rm den}+B_\eta
                                      =O_{\rm actual}(q).
\end{split}                                                    \tag{CRS15}
\]
For $B_\Gamma$, each argument $(s/2-1)/(q+r-j)$ has modulus
less than $1/2$, and $|\log(1+x)|\le2|x|$ on that interval.
For $B_v$, the intervals $[r+m,r+\Delta]$ overlap at most
$v+1$ times almost everywhere; integrate $2|\psi'(b/n_-)|$,
then substitute $t=b/n_-$. At $v=0$ the actual error is zero.
For $B_{\rm den}$, CRS6 yields, almost everywhere,
\[
 \left|f\left(\frac{r+u}{q+\ell-u}\right)
                    -f\left(\frac{r+u}{q}\right)\right|
 \le|\log(q/(q+\ell-u))|\le(\ell+\Delta)/Q .
\]
This includes the convergent limiting calculation at $r=0$.
The remaining regular difference is
$\int_0^\Delta[\eta((r+u)/q)-\eta(r/q)]\,du$.
For fixed $u$, the intervals $[r/q,(r+u)/q]$ overlap at most
$u+1$ times almost everywhere. Sum their absolute derivative
integrals and integrate in $u$ to obtain $B_\eta$.
These four calculations prove
\[
 \sum_{r=0}^M|h_r^{\rm prof}-c_r|\le B_{\rm sc}.        \tag{CRS16}
\]
Thus the original, possibly smaller finite radius satisfies
\[
\begin{gathered}
 E_c=E_1+\sum_{r=0}^M|h_r^{\rm prof}-c_r|
                 \le E_1+B_{\rm sc}=O_{\rm actual}(q),\\
 \boxed{\displaystyle\sum_{r=0}^M|\gamma_{r,s}-c_r|\le E_c.}
\end{gathered}                                                   \tag{CRS17}
\]
No independently known total or saturation argument enters this proof.

\subsection{Primitive and original endpoint weights}

Let $\mathfrak F(0)=0$, $\mathfrak F'=\psi$ be the original LRS15
primitive. Integration of CRS6 gives
\[
 \mathcal J(t):=\int_0^t f(u)\,du
 =4\mathfrak F(t)-2(1+t)\psi(t)+2a_0,\qquad
 2\mathcal J(1)=C_\partial.                            \tag{CRS18}
\]
Its original energy convention is retained:
\[
 \mathfrak F(t)=\frac{t^2}{4}F_{\rm energy}(2/t,\pi/t)
 +\frac{(1+t)^2}{2}\left(\frac32-\log(1+t)\right)-\frac34,
 \qquad F_{\rm energy}=-E_{\min}.
\]
For $R>0$, define
\[
\begin{split}
 A_\Delta(R)=\tfrac12[&(R+\Delta)^2\log(R+\Delta)
       -R(R+2\Delta)\log R\\
       &-\Delta^2\log\Delta-R\Delta],\qquad A_\Delta(0)=0 .
\end{split}                                                    \tag{CRS19}
\]
Direct differentiation gives
$A_\Delta'(R)=(R+\Delta)\log(1+\Delta/R)-\Delta$,
twice the correction subtracted in CRS13. Its limit at zero is zero.
Consequently
\[
 2\int_0^R c(x)\,dx=2\Delta q\mathcal J(R/q)-A_\Delta(R).
                                                               \tag{CRS20}
\]
The logarithmic part of CRS14 has derivative
$-\tfrac12\log(1+\Delta/x)$ and hence total variation
$\tfrac12[(q+\Delta)\log(q+\Delta)-q\log q-\Delta\log\Delta]$
on $[0,q]$. Its other term has variation at most $\Delta H_\eta$.
Therefore
\[
\begin{split}
 {\rm Var}_{[0,q]}c&\le V_c
 :=\tfrac12[(q+\Delta)\log(q+\Delta)-q\log q-\Delta\log\Delta]
                                    +\Delta H_\eta,\\
 V_c+|c_0|&=O(a\log q)=O(q).
\end{split}                                                    \tag{CRS21}
\]
For the size, the bracket equals
$q\log(1+\Delta/q)+\Delta\log(1+q/\Delta)$; use
$\log(1+x)\le x$. The profile is absolutely continuous including zero.

The original weights are
\[
 \mathcal T^{\le R}=\gamma_0+2\sum_{r=1}^R\gamma_r,\qquad
 \mathcal T^{>R}=2\sum_{r=R+1}^{q-1}\gamma_r+\gamma_q,\quad 0\le R<q .
                                                               \tag{CRS22}
\]
On an integer interval $[u,w]$, the doubled trapezoid
$c_u+2\sum_{u<r<w}c_r+c_w$ differs from $2\int_u^w c$
by at most $2{\rm Var}_{[u,w]}c$. On each unit interval the
absolute error is bounded by the integral of
$|c_j-c(x)|+|c_{j+1}-c(x)|$, at most twice its variation;
sum the unit intervals. The prefix in CRS22 has one extra
$c_R$ relative to this trapezoid and the tail has its negative.
Since $|c_R|\le|c_0|+V_c$, CRS17 proves simultaneously
\[
\begin{split}
 C^{\rm pre}(R)&=2\Delta q\mathcal J(R/q)-A_\Delta(R),\\
 C^{\rm tail}(R)&=2\Delta q[\mathcal J(1)-\mathcal J(R/q)]
                                    -A_\Delta(q)+A_\Delta(R),\\
 |\mathcal T^{\le R}-C^{\rm pre}(R)|,\
 |\mathcal T^{>R}-C^{\rm tail}(R)|
 &\le2E_c+|c_0|+3V_c=O_{\rm actual}(q).
\end{split}                                                    \tag{CRS23}
\]
The row weights cost at most $2E_c$. At $R=0$ the prefix is
$\gamma_0$ and its centre is zero, so the same bound holds directly.

\subsection{Moving-cutoff logarithms with a finite remainder}

Taylor's integral formula for $g(y)=(1+y)^2\log(1+y)$,
with $g'''(y)=2/(1+y)$, gives for every $R>0$
\[
\begin{split}
 A_\Delta(R)
 &=\tfrac{\Delta^2}{2}\log(R/\Delta)+\tfrac{3\Delta^2}{4}
  +\frac{R^2}{2}\int_0^{\Delta/R}
                         \frac{(\Delta/R-u)^2}{1+u}\,du,\\
 0&\le A_\Delta(R)-\tfrac{\Delta^2}{2}\log(R/\Delta)
                -\tfrac{3\Delta^2}{4}\le\frac{\Delta^3}{6R}.
\end{split}                                                    \tag{CRS24}
\]
This strengthens the received remainder and does not require $R\ge\Delta$.
For the prescribed $R_*=\lfloor a^{3/2}/\log a\rfloor$,
\[
\begin{gathered}
 \Delta^2/2=128q-1024a+1024,\\
 \log(R_*/\Delta)=\tfrac12\log a-\log\log a-\log16+O(1/a),\\
 \log(q/\Delta)=\log a-\log16+O(1/a),\qquad
 \Delta^3/R_*=O(a^{3/2}\log a)=O(q).
\end{gathered}
\]
The floor changes $\log R_*$ by $O(\log a/a^{3/2})$, absorbed
above. Inserting these expressions in CRS24 gives
\[
\begin{split}
 A_\Delta(R_*)&=64q\log a-128q\log\log a+O(q),\\
 A_\Delta(q)&=128q\log a+O(q),\\
 A_\Delta(q)-A_\Delta(R_*)&=64q\log a+128q\log\log a+O(q).
\end{split}                                                    \tag{CRS25}
\]
Thus
\[
 \boxed{\displaystyle
 \mathcal T^{>R_*}=2\Delta q\int_{R_*/q}^1 f(t)\,dt
   -64q\log a-128q\log\log a+O_{\rm actual}(q).}        \tag{CRS26}
\]
The entire equilibrium integral remains in this formula.
For fixed $1<\beta\le2$ and $R=\lfloor a^\beta\rfloor$,
the same computation gives
$A_\Delta(R)=128(\beta-1)q\log a+O_\beta(q)$.
The total is
$\Delta qC_\partial-128q\log a+O(q)
=16C_\partial aq-128q\log a+O(q)$, consistent with LRS24.

\subsection{Independent proper-source displacement}

Apply the proved profile at the independent original degree $k$,
its own centre $k/2-4$, order $s_0=1,k$, conductor and guards.
Put $q_k=(k+1)^2$, $b=8k+24$, $\Delta_k=16k-48$.
Here $c_r,E_c$ mean its degree-$k$ quantities through row $q_k+2b$.
ATA35--37 prove for the actual attained quotient $g(N)=\log\det Q_N$
that $0\le g(q_k+r)-g(q_k+r-1)\le\gamma_{r,s_0}$.
The exact four-cutoff displacement expands into negative intervals
$[0,b-1]$, $[1,b]$ and positive intervals
$[q_k,q_k+2b-1]$, $[q_k+1,q_k+2b]$. Keep each occurrence:
\[
\begin{split}
 L_k&=\sum_{r=0}^{b-1}c_r+\sum_{r=1}^{b}c_r,\qquad
 H_k=\sum_{r=q_k}^{q_k+2b-1}c_r+
                               \sum_{r=q_k+1}^{q_k+2b}c_r,\\
 -L_k-2E_c&\le\delta V_Q\le H_k+2E_c .
\end{split}                                                    \tag{CRS27}
\]
For the lower bound discard the nonnegative high increments and
bound the low increments by the native rows; CRS17 costs at most
$2E_c$. For the upper bound discard the low increments and apply
the same argument to the two high intervals. There is no additional
four-copy coefficient.

The low sum is the doubled trapezoid on $[0,b]$. CRS20--21 yield
$L_k=2\Delta_kq_k\mathcal J(b/q_k)-A_{\Delta_k}(b)+O(k\log k)$.
CRS6--7 give
$\mathcal J(t)=\tfrac t2(1-\log t)+O(t^{3/2})$ at zero.
Also $A_{\Delta_k}(b)=O(k^2)$: put $b=\Delta_kx$ in CRS19;
the $\log\Delta_k$ coefficients cancel and $x$ stays in a compact
positive interval. Hence
\[
\begin{split}
 L_k&=\Delta_kb[1+\log(q_k/b)]
        +O(\Delta_kb^{3/2}/\sqrt{q_k}+k^2)\\
    &=128q_k\log k+O(q_k).
\end{split}                                                    \tag{CRS28}
\]
Here $\Delta_kb=128k^2-1152=128q_k+O(k)$ and
$\log(q_k/b)=\log k-\log8+O(1/k)$.

For $q_k\le r\le q_k+2b$, CRS6 and CRS13 give
\[
 0\le\Delta_k f(1)-c_r
       \le2\Delta_k b/q_k+\Delta_k^2/(4q_k).
\]
Indeed $0\le f(1)-f(r/q_k)\le\log(r/q_k)\le2b/q_k$.
The subtracted correction equals
$\tfrac12\int_0^{\Delta_k}\log(1+u/r)\,du
\le\Delta_k^2/(4r)$. There are $4b$ high occurrences, so
\[
 H_k=4b\Delta_k f(1)+O(k)=512q_k f(1)+O(k).            \tag{CRS29}
\]
Combining CRS27--29 and $E_c=O_{\rm actual}(q_k)$ gives
\[
 \boxed{-128q_k\log k-O_{\rm actual}(q_k)
       \le\delta V_Q\le512q_k f(1)+O_{\rm actual}(q_k).}\tag{CRS30}
\]
This calculates the displacement allowances on the stated original
intervals. The independent proper-source standard-window determinant
remains in its original metric and is not evaluated by that displacement.

CRS25 distributes the existing native $-128q\log a$ between prefix
and tail. It still cancels the literal lower-root $+128q\log a$
in LRS22. Thus the existing complete-action identity remains
$J_a-\delta^{\rm ar}_{a,1}
=C_Bq^2+8q\log a+O_{\rm actual,h}(q)$.
The signed arithmetic correction, directional allocation exponents
and projected arithmetic responses acquire no assigned values from
this scalar calculation.
\endgroup
